\documentclass[11pt,letterpaper]{article}
\PassOptionsToPackage{hyphens}{url}
\usepackage[margin=.83in,top=.8in,bottom=.8in,headheight=15pt]{geometry}
\usepackage[T1]{fontenc}
\usepackage{lmodern,amsmath,amssymb,mathtools,bm,booktabs,array,graphicx,tikz}
\usetikzlibrary{arrows.meta,positioning,calc,decorations.pathreplacing,patterns}
\definecolor{ink}{HTML}{17344D}\definecolor{accent}{HTML}{087E8B}\definecolor{pale}{HTML}{EDF5F5}\definecolor{gold}{HTML}{AD6B19}
\providecommand{\dd}{\,\mathrm{d}}\providecommand{\R}{\mathbb{R}}
\providecommand{\E}{\mathbb{E}}\providecommand{\Var}{\operatorname{Var}}
\providecommand{\Cov}{\operatorname{Cov}}\providecommand{\diag}{\operatorname{diag}}
\providecommand{\norm}[1]{\left\lVert#1\right\rVert}
\providecommand{\abs}[1]{\left|#1\right|}
\usepackage{microtype,hyperref,fancyhdr,enumitem,amsthm,sectsty,needspace}
\allsectionsfont{\raggedright}
\hypersetup{colorlinks=true,linkcolor=ink,urlcolor=accent,pdfauthor={Zan Ahmad},pdftitle={The circulation / Review}}
\setlength{\parindent}{0pt}\setlength{\parskip}{6pt}\linespread{1.07}
\setlist{itemsep=3pt,topsep=4pt,leftmargin=*}
\emergencystretch=2em\clubpenalty=10000\widowpenalty=10000
\pagestyle{fancy}\fancyhf{}\fancyhead[L]{\small Mathematical Cardiology}\fancyhead[R]{\small Lecture 1}
\fancyfoot[L]{\footnotesize Zan Ahmad}\fancyfoot[R]{\footnotesize\thepage}\renewcommand{\headrulewidth}{.3pt}
\newtheorem{proposition}{Proposition}\newtheorem{theorem}{Theorem}\newtheorem{lemma}{Lemma}
\theoremstyle{definition}\newtheorem{definition}{Definition}\newtheorem{example}{Example}
\numberwithin{equation}{section}
\begin{document}
\section*{The circulation --- review sheet}
Use this after attempting the worked examples and practice. Explain each claim in words before calculating.

\section*{Key ideas}
\begin{itemize}
\item A circulation diagram states connections; it does not by itself specify geometry, pressures, or timing.
\item Arteries carry blood away from the heart; veins return it. Oxygen content does not define those names.
\item V is volume; Q is volume per time; P is pressure. Units distinguish the quantities.
\item A passive linear resistance obeys Q =\allowbreak{} (P\_in - P\_out)/R with R > 0.
\item Comparisons must state what is held fixed.
\item Input minus output changes stored volume; steady volume requires equal total flows.
\item A useful project question names an input, output, controlled comparison, validation check, and limitation.
\end{itemize}
\section*{Common errors and repairs}
\textbf{Use inlet pressure alone to calculate flow.} Subtract the outlet pressure and divide by the specified resistance.

\textbf{Interpret L/min as stored blood volume.} L/min is a rate. Multiply by time in minutes to obtain a passing volume in L.

\textbf{Call every oxygen-poor vessel a vein, or say venous blood is blue.} Use direction relative to the heart; pulmonary arteries are a counterexample. Blood is red at both oxygen levels.

\textbf{Say unequal instantaneous inflow and outflow violate conservation.} The difference changes stored volume; equality follows only when there is no accumulation.

\textbf{Treat the fixed-pressure resistor result as an exercise-response prediction.} The calculation holds boundary pressures fixed and omits pump dynamics and feedback.

\section*{Retrieve without your notes}
\begin{itemize}
\item Which additional number besides R is needed when you know only P\_in?
\item What unit should R have if P is in mmHg and Q in L/min?
\item What must be true of total inflow and outflow if stored volume is constant?
\end{itemize}
\section*{Readiness check}
\begin{itemize}
\item Trace the adult systemic and pulmonary circulation through the four heart chambers and identify the direction used to name arteries and veins.
\item Distinguish volume, volume flow rate, and pressure, giving a unit for each.
\item Calculate and interpret flow through one passive resistance using a pressure difference and stated assumptions.
\item Use an input-minus-output volume balance to test a proposed steady circulation.
\item State a small investigable project question with an input, an observable output, a controlled comparison, and a limitation.
\end{itemize}
\end{document}
